\documentclass[12pt,a4paper]{article}

\usepackage[margin=1in]{geometry}
\usepackage{amsmath,amssymb,mathtools}
\usepackage{array,booktabs}
\usepackage{enumitem}
\usepackage{listings}
\usepackage{xcolor}
\usepackage[hidelinks]{hyperref}
\usepackage{microtype}

\newcommand{\ket}[1]{\left|#1\right\rangle}

\lstdefinestyle{pythonstyle}{
    language=Python,
    basicstyle=\ttfamily\small,
    keywordstyle=\bfseries,
    commentstyle=\itshape,
    showstringspaces=false,
    breaklines=true,
    frame=single,
    columns=fullflexible
}
\lstset{style=pythonstyle}

\title{\textbf{Exercise Set: Simple Quantum Circuits with Qiskit}}
\author{}
\date{}

\begin{document}
\maketitle

\section*{Instructions}
This exercise set combines:
\begin{itemize}
    \item manual calculations of quantum states and gate operations,
    \item understanding and prediction of Qiskit programs,
    \item implementation of quantum gates and circuits using Qiskit,
    \item comparison of theoretical and simulated results.
\end{itemize}

Students should perform manual derivations before running the corresponding Qiskit programs wherever mentioned.

\section{Qubits and State Vectors}

\subsection*{Exercise 1: Basis states}
Write the vector representation of the following states manually:
\begin{enumerate}
    \item $\ket{0}$
    \item $\ket{1}$
    \item $\ket{00}$
    \item $\ket{01}$
    \item $\ket{10}$
    \item $\ket{11}$
\end{enumerate}

Using the tensor product, prove that
\[
\ket{01}=\ket{0}\otimes\ket{1}=\begin{pmatrix}0\\1\\0\\0\end{pmatrix}.
\]

\subsection*{Exercise 2: Normalization of a qubit}
Determine whether each state is a valid normalized qubit state.

\begin{enumerate}[label=(\alph*)]
    \item
    \[
    \ket{\psi}=
    \frac{1}{\sqrt{2}}\ket{0}+
    \frac{1}{\sqrt{2}}\ket{1}
    \]
    \item
    \[
    \ket{\psi}=
    \frac{1}{2}\ket{0}+
    \frac{\sqrt{3}}{2}\ket{1}
    \]
    \item
    \[
    \ket{\psi}=
    \frac{1}{2}\ket{0}+
    \frac{1}{2}\ket{1}
    \]
\end{enumerate}

Use
\[
|c_0|^2+|c_1|^2=1
\]
to justify your answer.

Then write Python code that verifies normalization numerically.

\subsection*{Exercise 3: Construct a state from $\theta$ and $\phi$}
Given
\[
\ket{\psi}=
\cos\left(\frac{\theta}{2}\right)\ket{0}
+
e^{i\phi}\sin\left(\frac{\theta}{2}\right)\ket{1},
\]
calculate the state manually for:
\begin{enumerate}
    \item $\theta=0,\ \phi=0$
    \item $\theta=\pi,\ \phi=0$
    \item $\theta=\frac{\pi}{2},\ \phi=0$
    \item $\theta=\frac{\pi}{2},\ \phi=\pi$
\end{enumerate}

Then implement:
\begin{lstlisting}
def state(theta, phi):
    ...
\end{lstlisting}

The function should return the two amplitudes.

\subsection*{Exercise 4: Measurement probabilities}
For
\[
\ket{\psi}=
\frac{\sqrt{3}}{2}\ket{0}
+
\frac{1}{2}\ket{1},
\]
calculate manually:
\[
P(0)
\qquad\text{and}\qquad
P(1).
\]

Then prepare the state in Qiskit and simulate 1000 measurements. Compare the experimental frequencies with the theoretical probabilities.

\section{Single-Qubit Gates}

\subsection*{Exercise 5: Pauli-X gate}
Given
\[
X=
\begin{bmatrix}
0&1\\
1&0
\end{bmatrix},
\]
calculate manually:
\[
X\ket{0}
\qquad\text{and}\qquad
X\ket{1}.
\]

Then implement both cases using Qiskit.

Complete:
\begin{lstlisting}
from qiskit import QuantumCircuit
from qiskit.quantum_info import Statevector

qc = QuantumCircuit(1)

# Add the required gate

state = Statevector.from_instruction(qc)
print(state)
\end{lstlisting}

\subsection*{Exercise 6: Two consecutive X gates}
Without running Qiskit first, predict the final state of
\[
X(X\ket{0}).
\]

Then verify using:
\begin{lstlisting}
qc.x(0)
qc.x(0)
\end{lstlisting}

Explain why
\[
X^2=I.
\]

\subsection*{Exercise 7: Hadamard gate}
Given
\[
H=
\frac{1}{\sqrt{2}}
\begin{bmatrix}
1&1\\
1&-1
\end{bmatrix},
\]
calculate manually:
\[
H\ket{0}
\qquad\text{and}\qquad
H\ket{1}.
\]

Then reproduce both cases in Qiskit and display the statevector.

\subsection*{Exercise 8: Two Hadamard gates}
Calculate manually:
\[
H(H\ket{0}).
\]

Then create a Qiskit circuit containing two consecutive Hadamard gates:
\[
\ket{0}\longrightarrow H\longrightarrow H.
\]

What is the final state? Show that
\[
H^2=I.
\]

\subsection*{Exercise 9: Pauli-Z gate}
Calculate:
\[
Z\ket{0}
\qquad\text{and}\qquad
Z\ket{1}
\]
for
\[
Z=
\begin{bmatrix}
1&0\\
0&-1
\end{bmatrix}.
\]

Then determine manually what happens in the circuit
\[
\ket{0}\xrightarrow{H}\xrightarrow{Z}.
\]

Verify your result using Qiskit.

\subsection*{Exercise 10: Phase versus measurement}
Consider:
\[
\ket{\psi_1}=
\frac{\ket{0}+\ket{1}}{\sqrt{2}}
\]
and
\[
\ket{\psi_2}=
\frac{\ket{0}-\ket{1}}{\sqrt{2}}.
\]

Answer:
\begin{enumerate}
    \item What is $P(0)$ for each state?
    \item What is $P(1)$ for each state?
    \item Are the two states identical?
    \item Why can computational-basis measurement not distinguish them directly?
\end{enumerate}

Create both states using Qiskit.

\section{Understanding Qiskit Code}

\subsection*{Exercise 11: Predict the output}
Study the following code without running it first:
\begin{lstlisting}
qc = QuantumCircuit(1)

qc.x(0)
qc.h(0)

state = Statevector.from_instruction(qc)

print(state)
\end{lstlisting}

Answer:
\begin{enumerate}
    \item What is the initial state?
    \item What is the state after \texttt{x(0)}?
    \item What is the final state after \texttt{h(0)}?
    \item What statevector do you expect?
\end{enumerate}

Then run the code and verify your prediction.

\subsection*{Exercise 12: Debug the circuit}
The intention is to prepare
\[
\frac{\ket{0}+\ket{1}}{\sqrt{2}}.
\]

A student writes:
\begin{lstlisting}
qc = QuantumCircuit(1)
qc.x(0)

state = Statevector.from_instruction(qc)
\end{lstlisting}

Identify the mistake and correct the code.

\subsection*{Exercise 13: Understand circuit ordering}
Consider:
\begin{lstlisting}
qc = QuantumCircuit(1)

qc.h(0)
qc.z(0)
qc.h(0)
\end{lstlisting}

First calculate the result manually. Then verify using Qiskit.

What familiar gate is equivalent to
\[
HZH?
\]

\section{Multiple Qubits and Tensor Products}

\subsection*{Exercise 14: Tensor-product calculation}
Calculate manually:
\[
\ket{0}\otimes\ket{1}.
\]

Use
\[
\ket{0}=
\begin{bmatrix}
1\\
0
\end{bmatrix},
\qquad
\ket{1}=
\begin{bmatrix}
0\\
1
\end{bmatrix}.
\]

Then verify using:
\begin{lstlisting}
np.kron()
\end{lstlisting}

\subsection*{Exercise 15: Product of superposition states}
Calculate:
\[
\ket{+}\otimes\ket{+},
\]
where
\[
\ket{+}=
\frac{\ket{0}+\ket{1}}{\sqrt{2}}.
\]

Express the result in the computational basis. Then implement a two-qubit Qiskit circuit that generates this state.

\subsection*{Exercise 16: Three-qubit superposition}
Start with
\[
\ket{000}.
\]

Apply a Hadamard gate to every qubit.

Determine manually the final state. Then implement:
\begin{lstlisting}
qc = QuantumCircuit(3)
\end{lstlisting}

How many computational-basis states have nonzero amplitude? What is the amplitude of each?

\section{Controlled Gates and Entanglement}

\subsection*{Exercise 17: CNOT truth table}
Construct the complete truth table for a CNOT gate. Assume the first qubit is the control and the second is the target.

Determine:
\[
\ket{00}\rightarrow ?,
\qquad
\ket{01}\rightarrow ?,
\qquad
\ket{10}\rightarrow ?,
\qquad
\ket{11}\rightarrow ?.
\]

Then verify each case using Qiskit.

\subsection*{Exercise 18: Create a Bell state}
Starting from
\[
\ket{00},
\]
apply:
\begin{enumerate}
    \item $H$ to qubit 0.
    \item CNOT with qubit 0 as control and qubit 1 as target.
\end{enumerate}

Calculate the state after each step. Show that the final state is
\[
\ket{\Phi^+}
=
\frac{\ket{00}+\ket{11}}{\sqrt{2}}.
\]

Implement the circuit in Qiskit and display both the circuit and the statevector.

\subsection*{Exercise 19: Bell-state measurement}
Modify the Bell-state circuit by adding measurements. Run it for:
\begin{itemize}
    \item 10 shots,
    \item 100 shots,
    \item 1000 shots,
    \item 10,000 shots.
\end{itemize}

Record the number of occurrences of \texttt{00} and \texttt{11}. Explain why \texttt{01} and \texttt{10} should ideally not occur.

\subsection*{Exercise 20: Generate another Bell state}
Design a circuit that generates
\[
\ket{\Psi^+}
=
\frac{\ket{01}+\ket{10}}{\sqrt{2}}.
\]

Tasks:
\begin{enumerate}
    \item Determine which additional $X$ gate is required.
    \item Draw the circuit manually.
    \item Implement it in Qiskit.
    \item Verify the statevector.
    \item Measure the circuit.
\end{enumerate}

\section{SWAP and Toffoli Gates}

\subsection*{Exercise 21: SWAP operation}
Prepare the state
\[
\ket{01}.
\]

Apply a SWAP gate. Predict the resulting state.

Then implement:
\begin{lstlisting}
qc.swap(0, 1)
\end{lstlisting}

Verify the result.

\subsection*{Exercise 22: SWAP decomposition}
A SWAP gate can be constructed using three CNOT gates:
\begin{lstlisting}
CNOT(q0,q1)
CNOT(q1,q0)
CNOT(q0,q1)
\end{lstlisting}

Test the circuit for:
\[
\ket{00},\quad \ket{01},\quad \ket{10},\quad \ket{11}.
\]

Compare the results with Qiskit's built-in:
\begin{lstlisting}
qc.swap(0, 1)
\end{lstlisting}

\subsection*{Exercise 23: Toffoli gate truth table}
The Toffoli gate has two control qubits and one target qubit.

Create its complete eight-row truth table for
\[
\ket{000}
\]
through
\[
\ket{111}.
\]

Identify the input combinations for which the target qubit changes.

Then verify using:
\begin{lstlisting}
qc.ccx(0, 1, 2)
\end{lstlisting}

\section{Measurements}

\subsection*{Exercise 24: Hadamard measurement}
Construct:
\begin{lstlisting}
qc = QuantumCircuit(1, 1)

qc.h(0)
qc.measure(0, 0)
\end{lstlisting}

Before running the simulator, predict the measurement distribution.

Then execute for 1000 shots and calculate:
\[
P(0)=\frac{N_0}{N},
\qquad
P(1)=\frac{N_1}{N}.
\]

Compare the results with theory.

\subsection*{Exercise 25: Deterministic versus probabilistic circuits}
Compare the following circuits.

\textbf{Circuit A}
\begin{lstlisting}
qc.x(0)
qc.measure(0, 0)
\end{lstlisting}

\textbf{Circuit B}
\begin{lstlisting}
qc.h(0)
qc.measure(0, 0)
\end{lstlisting}

For each circuit:
\begin{enumerate}
    \item Predict the outcome.
    \item Run 1000 shots.
    \item Explain why one circuit is deterministic and the other probabilistic.
\end{enumerate}

\section{Circuit Mapping}

\subsection*{Exercise 26: Quantum XOR circuit}
Classically,
\[
S=A\oplus B.
\]

Design a reversible quantum circuit using three qubits:
\[
\ket{A,B,0}
\]
such that the third qubit becomes
\[
A\oplus B.
\]

Tasks:
\begin{enumerate}
    \item Determine the required gates.
    \item Draw the circuit.
    \item Implement it in Qiskit.
    \item Test all four combinations of $A$ and $B$.
\end{enumerate}

\subsection*{Exercise 27: Quantum AND using Toffoli}
Use three qubits:
\[
\ket{A,B,0}.
\]

Design the circuit such that the third qubit becomes
\[
A\land B.
\]

Test:
\begin{center}
\begin{tabular}{ccc}
\toprule
$A$ & $B$ & Expected output \\
\midrule
0 & 0 & 0 \\
0 & 1 & 0 \\
1 & 0 & 0 \\
1 & 1 & 1 \\
\bottomrule
\end{tabular}
\end{center}

Use the Toffoli gate.

\section{Quantum Half-Adder}

\subsection*{Exercise 28: Manual half-adder design}
A classical half-adder produces
\[
S=A\oplus B
\]
and
\[
C=A\land B.
\]

Using four qubits
\[
\ket{A,B,S,C},
\]
where initially
\[
S=C=0,
\]
determine which quantum gates are required to calculate the sum and carry.

Do not use Qiskit initially. Draw the circuit manually.

\subsection*{Exercise 29: Implement the half-adder}
Implement the circuit from Exercise 28 using Qiskit.

Use:
\begin{lstlisting}
QuantumCircuit(4, 4)
\end{lstlisting}

Your circuit must use:
\begin{itemize}
    \item two CNOT gates for the sum,
    \item one Toffoli gate for the carry.
\end{itemize}

Test all four combinations:
\[
00,\quad 01,\quad 10,\quad 11.
\]

Create a table comparing theoretical and simulated results.

\subsection*{Exercise 30: Verify the half-adder equation}
For each input combination, verify:
\[
A+B=S+2C.
\]

Write Python code that automatically checks this equation for all four possible inputs. The program should print either \texttt{PASS} or \texttt{FAIL} for each case.

\section{Code Completion Exercises}

\subsection*{Exercise 31: Complete the Bell-state program}
Complete the missing statements:
\begin{lstlisting}
from qiskit import QuantumCircuit
from qiskit.quantum_info import Statevector

qc = QuantumCircuit(2)

# Step 1
_____________________

# Step 2
_____________________

state = Statevector.from_instruction(qc)

print(state)
display(qc.draw("mpl"))
\end{lstlisting}

The required output state is
\[
\frac{\ket{00}+\ket{11}}{\sqrt{2}}.
\]

\subsection*{Exercise 32: Complete the measurement program}
Complete the code so that a qubit is prepared in an equal superposition and measured 1000 times.

\begin{lstlisting}
from qiskit import QuantumCircuit
from qiskit_aer import AerSimulator

qc = QuantumCircuit(1, 1)

_____________________
_____________________

simulator = _____________________

result = simulator.run(
    qc,
    shots=_____________________
).result()

counts = _____________________

print(counts)
\end{lstlisting}

\section{Predict Before Running}

For each circuit below, do not run Qiskit first. Calculate the final state manually and then verify it.

\subsection*{Exercise 33}
\begin{lstlisting}
qc = QuantumCircuit(1)
qc.x(0)
qc.z(0)
\end{lstlisting}

\subsection*{Exercise 34}
\begin{lstlisting}
qc = QuantumCircuit(1)
qc.h(0)
qc.x(0)
\end{lstlisting}

\subsection*{Exercise 35}
\begin{lstlisting}
qc = QuantumCircuit(1)
qc.h(0)
qc.z(0)
qc.h(0)
\end{lstlisting}

\subsection*{Exercise 36}
\begin{lstlisting}
qc = QuantumCircuit(2)

qc.x(0)
qc.cx(0, 1)
\end{lstlisting}

Determine the final computational-basis state.

\subsection*{Exercise 37}
\begin{lstlisting}
qc = QuantumCircuit(2)

qc.h(0)
qc.cx(0, 1)
qc.x(1)
\end{lstlisting}

Find the final state algebraically and identify the corresponding Bell state.

\section{Challenge Problems}

\subsection*{Exercise 38: Create all four Bell states}
Create Qiskit circuits for:
\[
\ket{\Phi^+}=
\frac{\ket{00}+\ket{11}}{\sqrt{2}},
\]
\[
\ket{\Phi^-}=
\frac{\ket{00}-\ket{11}}{\sqrt{2}},
\]
\[
\ket{\Psi^+}=
\frac{\ket{01}+\ket{10}}{\sqrt{2}},
\]
and
\[
\ket{\Psi^-}=
\frac{\ket{01}-\ket{10}}{\sqrt{2}}.
\]

For each:
\begin{enumerate}
    \item Draw the circuit.
    \item Display the statevector.
    \item Measure 1000 shots.
    \item Compare the measurement results.
    \item Explain why computational-basis measurement cannot distinguish the $+$ and $-$ relative phases directly.
\end{enumerate}

\subsection*{Exercise 39: Gate matrices using Qiskit}
Use Qiskit to obtain the matrix representation of:
\begin{itemize}
    \item $X$,
    \item $H$,
    \item $Z$,
    \item CNOT.
\end{itemize}

For example:
\begin{lstlisting}
from qiskit.quantum_info import Operator
\end{lstlisting}

Compare each Qiskit-generated matrix with the corresponding matrix derived manually.

\subsection*{Exercise 40: Verify matrix multiplication against Qiskit}
For the circuit
\[
\ket{0}
\xrightarrow{H}
\xrightarrow{Z}
\xrightarrow{H},
\]
perform the calculation manually:
\[
HZH\ket{0}.
\]

Then implement the same circuit in Qiskit.

Compare:
\begin{enumerate}
    \item manual matrix multiplication,
    \item Qiskit statevector,
    \item final measurement result.
\end{enumerate}

\end{document}

